\documentclass[letterpaper,11pt]{article}
\usepackage[top=0.8in, bottom=0.8in, left=0.9in, right=0.9in]{geometry}
\usepackage[hidelinks]{hyperref}
\usepackage{lmodern}
\usepackage{parskip}
\usepackage{microtype}
\setlength{\parskip}{6pt}
\setlength{\parindent}{0pt}
\begin{document}
\begin{flushleft}
\textbf{\large Ajay Venkatesh} \\
Brooklyn, NY \\
+1 516 585 9013 \\
\href{mailto:av3855@nyu.edu}{av3855@nyu.edu}
\end{flushleft}
\vspace{10pt}
May 24, 2026
\vspace{6pt}

Hiring Manager \\
Paramount \\
New York, NY

\vspace{10pt}

Dear Hiring Manager,

My name is Ajay Venkatesh. I'm finishing my M.S. in Computer Engineering at NYU, with production software experience at PwC in Bengaluru, a summer SDE internship at Amazon, and ongoing on-campus work at NYU's Novel AI Technologies. I'm writing about the Software Engineer role on the Global Quality Engineering team.

The strongest match is the AI-tooling work I owned at Amazon. I deployed a \textbf{Retrieval-Augmented Generation (RAG)} pipeline with an \textbf{agentic workflow} that autonomously retrieved operational logs to reason over incident details, generating context-aware summaries and remediation steps. I implemented a \textbf{Model Context Protocol (MCP)} server for efficient tool-augmented retrieval, cutting AWS Athena query time and compute costs by \textbf{40\%}, and built an event-driven incident triage system on \textbf{AWS ECS Fargate, SQS, and SNS}. The shape of that work, AI-powered internal tools that make engineers faster on log analysis, summarization, and triage, is exactly what GQE is asking for.

In parallel at Campus Mesh I worked on \textbf{LLM-based retrieval pipelines (RAG)} for CampusIQ with \textbf{embedding-based semantic search} and structured data indexing, and shipped \textbf{Node.js microservices} and a \textbf{serverless backend} on AWS (Lambda, API Gateway, S3, CloudFront) across four environments with least-privilege IAM and KMS-managed keys. At PwC I owned production backend systems in \textbf{Java} on \textbf{Microsoft Azure} integrating with \textbf{Microsoft CRM} over \textbf{SQL Server}, delivered through code review and CI/CD with strict SDLC gates, which translates directly to the QE-integrated CI/CD environment this role lives in.

On overlap with your stack: \textbf{Java}, \textbf{Python}, \textbf{TypeScript}, plus \textbf{LLMs}, \textbf{embeddings}, \textbf{RAG} (LangChain, FAISS), \textbf{Hugging Face}, \textbf{Docker}, \textbf{Kubernetes}, and \textbf{AWS}, \textbf{Azure}, and \textbf{GCP}. AI-assisted development is how I work day to day; I use \textbf{Claude Code} and \textbf{Cursor} for refactors, debugging, and reasoning across unfamiliar code. \textbf{Kotlin} is the one stack item I haven't shipped in production yet, but I've picked up new backend languages quickly before.

What pulls me toward Paramount is the surface. Quality engineering at a global media company is one of the most underrated places where AI tooling can actually compound; better test generation, sharper log summarization, faster defect triage all show up downstream as cleaner releases and shorter incident times. Building that AI-powered tooling for an engineering organization that ships content to audiences worldwide, here in New York, is the kind of work I'd like to compound my career around.

I'd welcome the chance to discuss further. Thanks for considering my application.

\vspace{12pt}
Sincerely, \\
Ajay Venkatesh

\end{document}
